\section{Open challenges and research directions}
\label{sec:challenges}
The evidence map and the section-level synthesis point to six recurring
weaknesses in the current literature. These weaknesses cut across model
classes, integration patterns and application domains. They are not
about whether surrogate models can be useful; they concern how studies
are assembled, validated, compared and reported. The following
subsections focus on those gaps that appear repeatedly and that are most
relevant for future surrogate-assisted MES work.
\subsection{Decision-aware benchmarking}
\label{sec:ch:bench}
The most immediate gap is the lack of decision-aware benchmarks.
Predictive metrics are reported in many studies, but direct evidence on
decision quality is often limited to a few example runs. Speed-up claims
are common; regret distributions, feasibility rates and worst-case
violations are much rarer
\cite{Mylonopoulos202332697,Ruan2021221,Khaloie2025,Lim2025,Perera2019191}.
This makes it difficult to compare methods across papers, even when they
address similar tasks. A useful benchmark would fix the host problem,
the training set, the stress-test set and the solver-time budget, and
would require all methods to report the same decision-oriented outputs:
regret, feasibility rate, MIP-gap distribution, violation severity and
wall-clock savings. Recent learning-to-optimize work for
security-constrained dispatch and unit commitment shows that these
metrics can be reported in a systematic way
\cite{Wu2022231,Chen20244723,Chen2022,Wu2024391,Liang20246508}, but
the field still lacks a shared benchmark that would make results
directly comparable.
\subsection{Calibrated uncertainty for surrogate-assisted MOO}
\label{sec:ch:uq}
Uncertainty handling remains uneven. In surrogate-assisted MOO and in
chance-constrained formulations, the optimization logic depends on
posterior uncertainty, not only on the mean prediction. Yet calibration
is checked much less often than pointwise fit, and several studies
report breakdown precisely in those operating regions that are weakly
covered by the training set
\cite{Hu2020,Hu2021671,Volodina2022,Lawson201647,Hu2024,Buchanan2024,Subedi2026157}.
The field already has practical tools such as deep ensembles, Bayesian
last layers, conformal prediction and polynomial-chaos-based chance
constraints
\cite{Chen2023657,Liang20246508,Muhlpfordt20194806,Koirala20242284,Xu20212696}.
The unresolved issue is how to combine these tools with surrogate-
assisted search and uncertainty-constrained optimization so that the
reported uncertainty still reflects what the optimizer sees at the
relevant points in the design or operating space.
\subsection{Out-of-distribution generalisation across capacity mixes}
\label{sec:ch:ood}
Planning studies use surrogates to say something about systems that do
not yet exist in their final form. This means that the relevant test is
often out-of-distribution by construction: another capacity mix, another
weather year, another operating regime, or another policy setting. Many
surrogates, however, are trained on historical operation or on samples
generated from a much narrower design space
\cite{Park2019,Reynolds2017704,Subedi2026157,Khaloie2025,Lim2025,Mylonopoulos202332697}.
Two practical responses appear in the literature. One is to build more
physics-aware surrogates that preserve balances, monotonicity or other
structural constraints
\cite{Park2019,Su2022315,Liu20223726,Petrova2025,Du2025}. The other
is to fine-tune a base surrogate with a small number of new
high-fidelity runs when climate, year or capacity mix changes
\cite{Chen202514792,De Castro20247710}. What is
still missing is a common way of reporting this distinction, so that
within-distribution accuracy is not confused with planning-relevant
robustness.
\subsection{Standardised public datasets and reproducibility}
\label{sec:ch:data}
The reproducibility audit in Section 8.6 showed a familiar pattern:
many studies describe the case study and the model idea, but not the
full artefacts needed to reproduce the result. Our evidence audit found
few public datasets, especially for district heating, sector-coupled MES
and other settings in which surrogate-assisted optimization would be
most useful.
The bulk of MES surrogate work is still built around individual case
studies. A useful minimum standard would be a public optimization
instance, the corresponding training data, an explicit train/validation/
test split, solver settings and a small number of reference truth-model
runs. Existing load and renewable forecasting datasets
\cite{Eseye201991463,Faraji2020157284,Hanifi2024,Sideratos2020,Tian2020,Madhukumar20228891}
show that this level of release is feasible.
\subsection{MCDM under surrogate uncertainty}
\label{sec:ch:mcdm}
The literature on multi-criteria decision making after surrogate-
assisted MOO is still thin, although this is exactly the step at which a
planner or designer needs a final decision. Some studies combine
response-surface MOO with a subsequent multi-criteria selection
\cite{Roy2020}, but the reviewed literature rarely propagates surrogate
uncertainty into that final ranking. In most multi-objective studies,
the Pareto set is presented and a compromise design is selected without
quantifying whether surrogate error could change the ordering.
For MES planning, this is a concrete gap: uncertainty should not stop at
the Pareto front but should be carried through to the final ranking or
selection step.
\subsection{Open-source libraries and toolchain integration}
\label{sec:ch:tooling}
The final gap is practical rather than conceptual. Several strands of
the literature already have working toolchains -- for example
polynomial-chaos OPF, learning-to-optimize for dispatch, or surrogate-
assisted evolutionary design -- but each strand tends to build its own
data pipeline, evaluation logic and reporting format
\cite{Muhlpfordt20194806,Mitrovic2023__2,Mitrovic2023__3,Wu2022231,Chen20244723,Chen2022,Hirvonen2017323,Schulte2020}.
This fragmentation makes reuse difficult and slows down comparison
across model classes and applications. A shared open-source layer would
not need to replace existing host optimizers; it would be enough to
standardize surrogate input/output interfaces, validation routines and
reporting formats. Frameworks such as PyPSA, Calliope, oemof, SpineOpt
and Backbone already provide a basis on the optimization side. The
surrogate side is still scattered.
